\documentclass[12pt]{article}
\author{Hans Halvorson}
\title{Niels Bohrs Plads i den Filosofiske Tradition}
\date{September 15, 2022}
\usepackage{hyperref}

\begin{document}

\maketitle

Alle akademikere ved, at Niels Bohr var en stor fysiker. I dag vil jeg
dog tale om Niels Bohr som filosof. Selvom Bohr ikke ville have kaldt
sig selv for en filosof, er det uomtvisteligt, at han passer ind i en
tradition af filosofiske videnskabsmænd, som inkluderer René
Descartes, Isaac Newton, Gottfried Leibniz, og H.C. Ørsted. Faktisk
kan vi sige det endnu stærkere: Bohr var en ganske original filosof og
han beskæftigede sig med de store spørgsmål om menneskelig erkendelse
som havde optaget de største europæiske filosoffer såsom John Locke,
David Hume, og Immanuel Kant.

Jeg vil beskrive Bohrs bidrag på baggrund af G.F.W. Hegels beretning om
filosofiens historie, som blev skrevet cirka hundrede år før Bohr indtog
scenen. Ifølge Hegel har filosofiens og videnskabens historie udgjort en
kamp for at opnå objektiv viden om virkeligheden. Hegel mente, at
tidligere filosoffer såsom Descartes og Spinoza bidrog til opnåelsen af
dette mål, især med deres ide om, at videnskabens mål er at beskrive
virkeligheden \emph{sub specie aeternitatis}, dvs. under evighedens
synspunkt, eller som man nu siger, fra et ''god's eye view''. Ifølge
Hegel mistede andre filosoffer modet og overgav sig til subjektiviteten.
Herunder nævnede han Hume, Kant, og sin umiddelbare forgænger, Johann
Gottlieb Fichte.

Jeg er totalt uenig med Hegel i, at idéhistorien er en proces,
hvorigennem ''den absolutte ånd'' (\emph{Geist}) realiserer objektiv
viden om sig selv. Jeg er dog enig med Hegel i, at der findes en
vedvarende dialektik mellem objektive og subjektive elementer i viden,
og at denne dialektik udspiller sig i videnskabens historie. Det er
netop her, jeg mener, at Bohr ydede sit største bidrag. Ligesom Hegel
var Bohr en fortaler for videnskab og for objektiv viden. Desuden var
Bohr af den opfattelse, at fysik er den videnskab, som er bedste egnet
til at give objektiv viden. Men Bohrs syn på objektiv viden er ikke det
samme som Hegels. Bohrs erkendelsesteori er snarere en udvikling af
ideer han arvede fra sine danske forgængere.

\hypertarget{metode}{%
\section{Metode}\label{metode}}

Jeg vil argumentere for, at Bohr blev påvirket af sine danske
forgængere, i særdeleshed Søren Kierkegaard, Rasmus Nielsen, og Harald
Høffding. Men det betyder, at jeg stiller mig selv mod Bohrs mest
fremragende filosofiske eksponent, David Favrholdt. I sit skrift ''Niels
Bohr's Philosophical Background'' fra 1992 og igen i sin bog
\emph{Filosoffen Niels Bohr} fra 2009 argumenterede Favrholdt for, at
Bohr stort set intet skylder tidligere filosoffer. Favrholdt hævdede
nemlig, at hverken Kierkegaards eller Høffdings ideer skulle have
påvirket Bohr overhovedet.

Hvad der er interessant, er Favrholdts metode: for at underbygge, at
Kierkegaards ideer ikke påvirkede Bohr, mente Favrholdt, at det var nok
at vise, at Bohr ikke læste mange af Kierkegaards skrifter og kun
sjældent refererede til Kierkegaard. Og Favrholdt kendte de relevante
kilder bedre end nogen anden: hans argumenter er ret overbevisende,
nemlig at der er få spor, som etablerer et direkte forhold mellem
Kierkegaard og Bohr. Derudover gav Favrholdt belæg for, at Bohr ikke
havde et særlig tæt forhold til Høffding, i det mindste ikke efter at
Bohr tog Høffdings Filosofikum.

Jeg synes, at Favrholdts tilgang er interessant, men at han var for
optaget af de materielle forbindelser fremfor idémæssige ligheder. Lad
mig give en analogi. Antag at du vil afgøre, om en far har påvirket sin
søn. Ifølge Favrholdts tilgang kan du spørge sønnen direkte: ''har din
far påvirket dig?'' Antag at sønnen svarer, ''nej, jeg hader min far og
jeg ville aldrig lade mig påvirke af ham.'' Antag videre, at sønnen
aldrig lyttede til sin far og flyttede hjemmefra, da han var ung. I så
fald må du konkludere -\/- hvis du tilgår spørgsmålet som Favrholdt -\/-
at sønnen ikke blev påvirket af sin far.

Der er dog en anden slags evidens, som peger på, at sønnen ikke desto
mindre blev påvirket af sin far. For eksempel hvis sønnens måde at tale
på minder om farens eller hvis sønnen har lignende præferencer, er det
underordnet, at sønnen benægter farens påvirkning eller ikke tilbragte
meget tid med sin far. Selv hvis de fysiske forbindelser er tynde, kan
det være interessant at tænke over de træk, som faren og sønnen har
tilfælles. På en lignende måde kan det være interessant at tænke over de
idemæssige træk, som Bohr har tilfælles med tidligere filosoffer. Derfor
leder jeg ikke efter konkrete spor, som beviser, at Bohr læste eller
hørte noget, som vi ikke var opmærksomme på. Jeg vil hellere gå direkte
til de kendte kilder for at sammenligne ideer.

\hypertarget{bohrs-centrale-iduxe9}{%
  \section{Bohrs centrale idé}\label{bohrs-centrale-iduxe9}}

Hvis Bohr senere i livet var blevet spurgt om, hvad hans vigtigste
filosofiske idé var, så ville han helt sikkert have sagt
''komplementaritet'' -\/- et ord som han begyndte at bruge i 1920-erne
for at optegne, at man ofte har brug for flere forskellige beskrivelser
for at indfange en situations helhed. Jeg er dog ikke sikker på, at Bohr
havde ret om sig selv. Jeg er enig med Aage Petersen og Léon Rosenfeld,
om at komplementaritet ikke skal ses som en lære, man kan indfange med
få ord, men snarere som en generel holdning til hvordan man bør bedrive
videnskab. Så jeg ville ikke kalde komplementaritet for en filosofisk
idé -\/- selvom jeg ville argumentere for, at komplementaritet viser sig
i flere af Bohrs filosofiske ideer.

Det kunne være givtigt at gå længere tilbage i Bohrs egen udvikling som
en filosofisk tænker -\/- da han begyndte at tale om ''kvantefysikkens
erkendelsesmæssige lære''. Denne lære handler om, at mennesket ikke kan
opnå objektiv viden i traditionel forstand, og Bohr fremfører det på
denne måde:

\begin{enumerate}
\def\labelenumi{\arabic{enumi}.}
\item
  Objektiv viden i traditionel forstand forudsætter, at det vidende
  subjekt er fuldstændigt adskilt fra den virkelighed (dvs. objektet),
  som det betragter.
\item
  (Kvantepostulatet) Virkningsenheden har en mindste mulig værdi, som er
  større end nul. Denne værdi kaldes for virkningskvantet.
\item
  En fuldstændig adskillelse af to fysiske objekter forudsætter, at
  virkningsenheden kan betragtes som vilkårlig lille.
\item
  Konklusion: Det er umuligt at nå objektiv viden (i traditionel
  forstand).
\end{enumerate}

Hvis Bohr stoppede her, kunne han kaldes for en skeptiker, dvs. én som
ikke tror på muligheden for objektiv viden. Men Bohr stoppede ikke med
det negative resultat. Han indså, at subjektet kan give afkald på
fuldstændighed til gavn for nøjagtighed. For at være mere præcis kan
subjektet opnå en entydig beskrivelse af ét aspekt af situationen, hvis
det opgiver at beskrive alle aspekter samtidigt.

Her følger jeg David Favrholdt og Don Howard, som har gengivet Bohrs
ræsonnement på dette punkt således:

\begin{enumerate}
\def\labelenumi{\arabic{enumi}.}
\item
  Éntydig beskrivelse forudsætter et skarpt skel mellem meddelelsen og
  personen, som giver den.
\item
  Man kan sætte et sådant skel ved at specificere konteksten, hvori
  meddelelsen formuleres.
\end{enumerate}

Begge pointer kræver nærmere forklaring. Den første pointe synes måske
kontroversiel, men den er ikke ejendommelig for Bohr. Faktisk bruger
engelsktalende filosoffer frasen ''propositional knowledge'' for at
antyde, at det der vides objektivt (the proposition), kan adskilles fra
den vidende. Desuden virker det intuitivt, at en meddelelse kun er
éntydig, hvis den er frigjort fra skjulte betydninger eller
konnotationer, som kan findes i meddelerens hoved.

Det første punkt leder umiddelbart til det andet. En meddeler
fremstiller en form for artefakt (f.eks. en skreven sætning), som han
eller hun håber, at modtageren kan udlede den tilsigtede betydning fra.
I mange tilfælde er den pågældende sætning ikke i sig selv nok til at
genskabe betydningen fra: modtageren skal også kende de betingelser,
under hvilke sætningen blev fremstillet. Du kan for eksempel kun udlede
betydningen af sætningen ''skattekisten ligger to kilometer øst for
mig'', hvis du ved, hvem der siger den og hvor han eller hun befinder
sig.

Bohr var dybt opmærksom på disse pointer om betydningsfuldt sprogbrug og
tog det skridtet videre. Han mente, at når fysikere samarbejder og taler
om eksperimenternes resultater, når deres kommunikation kun i mål, hvis
de har en fælles forståelse af de eksperimentelle betingelser. Det er
altså denne specifikation af betingelserne, som muliggør entydigt
sprogbrug.

Disse pointer kan virke ret banale, men de udgør et stridspunkt blandt
filosoffer. Rationalister såsom Spinoza og Hegel mente, at videnskabens
mål er ''absolut viden'', dvs. viden der ikke er betinget af den
videndes tilfældige tilstand eller kontekst. Denne opfattelse af
videnskabens mål blev afvist af Møller, Kierkegaard, Nielsen, og
Høffding, og Bohr ser ud til at følge trop.

Så vidt jeg ved, omtalte Bohr aldrig Hegel. Ikke desto mindre kan det
siges, at Bohrs opfattelse af objektiv viden stiller ham på Kierkegaards
side i debatten med Hegel. Bohr tog også de Kierkegaardske ideer
skridtet videre. For det første er Bohrs kritik af den Hegelske
opfattelse baseret på nye opdagelser i fysik i stedet for vage
filosofiske overvejelser. For det andet giver Bohr os ikke blot en
negativ kritik -\/- han udvikler også et positivt syn på, hvordan
mennesket kan opnå objektiv viden, selvom det ikke er i stand til at
opnå Hegels absolutte viden.

\hypertarget{fra-sibbern-til-kierkegaard}{%
\section{Fra Sibbern til
Kierkegaard}\label{fra-sibbern-til-kierkegaard}}

I Kierkegaardforskningen uden for de nordiske lande bliver hans
forgængere sjældent nævnt ved navn, hvilket giver indtrykket af, at
Kierkegaards ideer opstod ud af det blå. I virkeligheden arvede
Kierkegaard en filosofisk tradition, som inkluderer Kierkegaards lærere
Sibbern og Møller. Sibbern og Møller havde blandt andet tilfælles, at de
begge afviste Hegels opfattelse af objektiv viden.

Møller døde ung og udviklede aldrig en erkendelsesteori, som kunne
erstatte Hegels. Til gengæld levede Sibbern længe og han udviklede en
enormt kompleks filosofi. Desværre er jeg på nuværende tidspunkt ikke i
stand til at komme ind på et eventuelt forhold mellem Sibberns ideer og
Bohrs. Sibberns filosofi er sært vanskelig og ikke særlig videnskabelig.
Jeg springer ham derfor over for nu og fortsætter med Kierkegaard -- og
hermed vender jeg mig væk fra Favrholdt.

Hvor Favrholdt reducerer forholdet mellem Kierkegaard og Bohr til, at
''de begge var danskere'', mener jeg, at et nærmere studie af deres
værker viser mange ligheder i deres tankegange. Der er ikke tid til at
gå grundigt igennem Kierkegaard og Bohrs ideer i dag, og jeg vil derfor
nøjes med at belyse et par lighedspunkter. Desuden vil jeg komme ind på,
hvordan disse punkter går igen i Nielsen og Høffding.

Hegel tog den traditionelle opfattelse af objektiv viden til dens
logiske yderlighed: subjektet skal opnå et ''god's eye view'' eller
''view from nowhere'', dvs. det skal se virkeligheden, som den er i
sig selv og ikke som den ser ud fra et eller andet perspektiv. Desuden
spiller subjektets præferencer, følelser, køn, nationalitet,
osv. ingen rolle ifølge Hegel, eftersom disse faktorer er irrelevante
for videns indhold.

Eller det er i det mindste sådan Møller, Sibbern, og Kierkegaard forstod
Hegels ideal og de konkluderede, at dette hverken er opnåeligt eller
rimeligt. Faktisk hævdede Kierkegaard, at viden forudsætter, at
mennesket ikke er en passiv tilskuer til tilværelsens drama, sådan som
Hegel forestillede sig det. For eksempel argumenterede Kierkegaard for,
at en person kun kan få sikker viden om sit forhold til en anden, hvis
han eller hun elsker den anden.

Kierkegaard imødegik effektivt Hegels ideal ved at opridse
subjektivitetens betydning for livet som menneske. Det er velkendt, at
Kierkegaard favoriserede humanistiske områder (poesi, kærlighed,
familie, religion, osv.) over videnskaben. Det er derfor nemt at forstå,
hvorfor Favrholdt afviste, at der er fælles træk i Kierkegaard og Bohrs
tankegange. Favrholdt kunne ikke se videre end, at Kierkegaard var
videnskabens modstander mens Bohr var videnskabens forkæmper. Hvad
Favrholdt ikke indså var, at Rasmus Nielsens videreudvikling af
Kierkegaards filosofi inkluderede et positivt videnskabssyn.

\hypertarget{fra-kierkegaard-til-huxf8ffding}{%
  \section{Fra Kierkegaard til
    Høffding}\label{fra-kierkegaard-til-huxf8ffding}}

Mellemleddene mellem Bohr og Kierkegaard udgør den store ufortalte
historie om Bohrs plads i den filosofiske tradition. Takket være Jan
Faye har vi et overblik over Bohrs og Høffdings forhold. Jeg vil dog
ikke blande mig i Fayes og Favrholdts debat, om hvorvidt forholdet var
tæt nok til, at man med rette kan tale om et påvirkningsforhold. (Faye
bekræfter mens Favrholdt benægter.) Ikke desto mindre mener jeg, at Faye
har givet rigeligt bevis for, at den ideverden som Bohr levede i til
dels blev skabt af Høffding og derfor også af dem, som påvirkede
Høffding. Her skal ikke kun nævnes Kierkegaard og hans forgængere men
også det forsømte led mellem Kierkegaard og Høffding: Rasmus Nielsen.

Nielsen var Kierkegaards samtidig på Københavns Universitet, men i
modsætning til den nu verdensberømte forfatter blev Nielsen professor på
Universitetet---en stilling han besad i fire årtier. I løbet af sin
første år som professor blev Nielsen påvirket af Kierkegaard og det fik
ham til at gøre op med Hegels filosofi og udvikle en
Kierkegaard-inspireret videnskabsteori.

Kierkegaard selv havde ringe interesse for dagens videnskab og Hegel var
stort set uvidende om nye udviklinger i naturvidenskaben. Det var
Nielsen til gengæld ikke. Han forblev interesseret i fagvidenskaberne og
han skrev flere bøger om fysik, kemi, matematik, osv. Derudover
forelæste Nielsen i mange år på det indledende Filosofikum, hvor han
kunne regne blandt sine studerende Christian Bohr (Niels' far) og Harald
Høffding.

Lad os kigge på, hvordan Nielsen anvendte Kierkegaardske ideer i sin
videnskabsteori.

Rationalistiske filosoffer såsom den tyske Christian Wolff antog, at Gud
udstyrede mennesket med begreber, som svarer til virkelighedens
grundlæggende struktur. Denne antagelse kaldte Immanuel Kant for
''dogmatisme'' fordi Gud -- forudsat at han overhovedet eksisterer --
ikke har givet mennesket en garanti for, at det vil kunne forstå
virkeligheden, som den er i sig selv. Men Kant var heller ikke glad for
den empiriske opfattelse, efter hvilken mennesket kun har et begreb,
hvis det svarer til en sansefornemmelse. Empiristen David Hume havde
argumenteret overbevisende for, at man ikke har sansefornemmelser af
årsag og substans, og han konkluderede, at objektiv viden er umulig.
Derimod tog Kant for givet, at fysikere som Isaac Newton havde opnået
objektiv viden. Så Kant stillede sig selv den opgave at forklare
muligheden for objektiv viden.

Immanuel Kants løsning var, at objektiv viden muliggøres af en
begrænsningsstrategi: man begrænser sig til den fænomenale verden, dvs.
den verden som bliver ''objektiveret'' (dvs. gjort til et vidensobjekt)
af det menneskelige sind. Man bruger ofte her en analogi: verden i sig
selv er som en formløs klump kagedej og det menneskelige sind er som en
kageudstikker, der påtvinger dejen en bestemt form. Derfor er menneskets
''former for forståelse'' med til at skabe den objektive verden.

Flere filosoffer efter Kant syntes ikke, at hans syn var en løsning på
problemet om videnskabens mulighed, men snarere en begrænsning af
videnskabens domæne. De ville bevise, at objektiv viden ikke er
begrænset til den fænomenale verden, men kan ''uendeliggøres'' ved hjælp
af en ''uendelig objektiveringsproces''. Det var i hvert fald Hegels
opfattelse. Derudover mente Hegel, at han havde fuldført processen selv
og derved havde opnået absolut viden.

De danske filosoffer -- Møller, Sibbern, og Kierkegaard --
latterliggjorde Hegels indbildskhed. Men betyder det også, at de
benægtede muligheden for objektiv viden? Det gjorde Rasmus Nielsen ikke.
Han videreførte Kants ide om, at videnskab muliggøres af, at mennesket
foretager en objektiveringsproces. Dog er al menneskelig objektivering
delvis. Kun Gud kan fuldføre objektiveringsprocessen. For Gud er
virkelighed og forestilling én og den samme.

Nielsens opfattelse af objektivering er tydeligt beskrevet i Høffdings
filosofihistorie. Det er interessant, at Høffding brugte anledningen til
at fremføre sin egen kritik af Nielsen. Kritikken er dog langt fra et
forsvar af Hegels opfattelse. I stedet for at forsvare ideen om, at
mennesket kan nå ''god's eye view'', argumenterede Høffding for, at der
ikke findes et ''god's eye view''. Han hævdede nemlig, at ideen er
meningsløs.

Selvom Favrholdt benægtede, at hverken Kierkegaard eller Høffding havde
en væsentlig indflydelse på Bohr, brugte han et helt kapitel af
\emph{Filosoffen Niels Bohr} til at forklare, at Bohr afviste ''the
god's eye view''. Desuden hævdede Favrholdt, at denne afvisning var
nøglepunktet i Bohrs langvarige debat med Einstein. Men hvordan, må man
spørge, var Bohr dristig nok til at indtage denne holdning, som kunne
forekomme at være i modstrid med videnskabens ånd? Jeg påstår, at Bohr
så Høffding og indirekte Nielsen som fortalere for en form for
objektivitet, der adskiller sig fra Hegels absolutte viden, men er værd
at stræbe efter.

\hypertarget{hvad-suxe5}{%
\section{Hvad så?}\label{hvad-suxe5}}

Jeg er klar over, at denne slags abstrakte filosofiske overvejelser ikke
falder i enhver smag. Kan videnskaben ikke bare fortsætte, uden at man
skal snakke om Hegel, ''god's eye view'', eller dialektikken mellem
subjektivitet og objektivitet?

Det kan den for det meste. I perioderne hvor videnskab forandrer sig,
opstår der dog spørgsmål af en mere filosofisk karakter. Det var
tilfældet med kvantemekanikkens udvikling og det er tilfældet igen den
dag i dag.

Einstein sagde tit -- og mest markant i den berømte
Einstein-Podolsky-Rosen artikel fra 1935 -- at kvantemekanikken ikke
var en fuldkommen teori. Dvs., at kvantemekanikken var en mangelfuld
teori, som burde erstattes. For at give et konkret eksempel: hvis
Einstein var formand for Grundforskningsfonden, tror jeg ikke, at han
ville investere i kvanteinformationsteori eller strengteori eller
kvanteoptik, fordi han ville have været af den opfattelse, at de alle
sammen bygger på et forkert grundlag. Han ville sandsynligvis have
foretrukket at støtte udviklingen af alternativer til
kvantemekanikken, såsom David Bohms bølgeteori eller Ghirardi, Rimini,
og Webers ''spontaneous collapse" teori.

Hvorfor kunne Einstein ikke acceptere kvanteteorien? Her tilslutter jeg
mig Favrholdt: Einstein ville ikke opgive søgen efter ''god's eye view''
og han vidste godt, at kvantemekanikken ikke giver det. Det forklarer
også, hvorfor stadig flere nuværende filosoffer og fysikere er blevet
mistroiske over for kvantemekanikken. Den trang, som var udgangspunktet
for rationalisternes filosofiske indsats, er fortsat operationel -\/- og
især blandt dem, som ikke kender til filosofiens historie.

Vi ved, hvilke slags bevillinger Niels Bohr ville have godkendt, fordi
vi kan kigge på, hvilke projekter han støttede som institutleder. Mit
beskedne bud er, at Bohrs succes som institutleder til dels skyldtes, at
han tog beslutninger på baggrund af en nøje gennemtænkt filosofi.

\end{document}

%%% Local Variables:
%%% mode: latex
%%% TeX-master: t
%%% End:
